\chapter{Mutation Testing cu PITest}
\label{ch:pitest}

\section{Cum rulăm PITest}

Spre deosebire de JaCoCo, PITest nu rulează în faza standard \texttt{test} a Maven. Trebuie invocat explicit:

\begin{lstlisting}[style=plainstyle]
mvn test-compile org.pitest:pitest-maven:mutationCoverage
\end{lstlisting}

După rulare, raportul HTML se găsește în \texttt{target/pit-reports/index.html}. PITest generează în plus un raport XML, util pentru integrarea în CI/CD.

Pentru clasa noastră, durata totală a unei rulări PITest e între 5 și 15 minute, în funcție de încărcarea CPU-ului. Procesul intern e următorul:

\begin{enumerate}
    \item PITest analizează clasa \texttt{WalletServiceImpl} și identifică toate locurile unde poate aplica un mutator (loop-uri, condiții, calcule, return-uri etc.).
    \item Generează variante mutate ale codului sursă (211 mutanți, în cazul nostru).
    \item Pentru fiecare mutant, recompilează clasa și rulează suita de teste relevantă.
    \item Pentru fiecare test, înregistrează rezultatul: \textit{KILLED} (testul a căzut $\Rightarrow$ mutantul a fost detectat), \textit{SURVIVED} (toate testele au trecut $\Rightarrow$ mutantul a scăpat), \textit{NO\_COVERAGE} (mutantul nu e atins de niciun test) sau \textit{TIMED\_OUT}.
    \item Generează raportul HTML cu evidențierea liniilor pe culori în funcție de starea mutanților.
\end{enumerate}

\section{Rezultate generale}

Raportul final pentru \texttt{WalletServiceImpl} arată:

\begin{table}[!htb]
\centering
\caption{Rezultate generale PITest pe \texttt{WalletServiceImpl}.}
\label{tab:pitest-summary}
\begin{tabular}{lc}
\toprule
\textbf{Metric} & \textbf{Valoare} \\
\midrule
Mutanți generați & 211 \\
Mutanți omorâți & \textbf{125} \\
Mutanți supraviețuitori & 86 \\
Line Coverage (PITest) & 99\% (177/179) \\
\textbf{Mutation Score} & \textbf{60\%} (125/211) \\
Test Strength & 70\% (146/210) \\
\bottomrule
\end{tabular}
\end{table}

\begin{figure}[!htb]
    \centering
    \fbox{\includegraphics[width=0.9\linewidth]{Pit_Test.png}}
    \caption{Captură din raportul PITest HTML, cu sumarul general.}
    \label{fig:pitest-screenshot}
\end{figure}

Distincția între \textit{Mutation Score} și \textit{Test Strength}:

\begin{itemize}
    \item \textbf{Mutation Score} = \texttt{killed / total\_generated}. Include și mutanții care nu sunt atinși de niciun test.
    \item \textbf{Test Strength} = \texttt{killed / total\_covered\_by\_tests}. Doar pentru mutanții pe care testele îi ating.
\end{itemize}

Diferența între cei doi indici (60\% vs. 70\%) ne spune că aproape toate liniile sunt atinse de cel puțin un test (Line Coverage 99\%), dar testele nu sunt suficient de \enquote{stricte} ca să detecteze toate alterațiile.

\section{Defalcare pe tipuri de mutatori}

PITest grupează mutanții după tipul mutatorului. Tabelul \ref{tab:mutators} prezintă defalcarea, iar figura \ref{fig:mutator-distribution} oferă reprezentarea grafică.

\begin{table}[!htb]
\centering
\caption{Distribuția mutanților pe tipuri de mutatori.}
\label{tab:mutators}
\begin{tabularx}{\linewidth}{Xccc>{\centering\arraybackslash}p{2cm}}
\toprule
\textbf{Mutator} & \textbf{Generat} & \textbf{Omorât} & \textbf{Supr.} & \textbf{Kill rate} \\
\midrule
\textsc{NullReturnVals}             & 18 & 17 & 1  & 94\% \\
\textsc{VoidMethodCall}             & 11 & 11 & 0  & \textbf{100\%} \\
\textsc{RemoveCondEqual}            & 53 & 30 & 22 & 57\% \\
\textsc{RemoveCondOrder}            & 21 & 12 & 8  & 57\% \\
\textsc{ConditionalsBoundary}       & 21 & 7  & 13 & 33\% \\
\textsc{EmptyObjectReturn}          & 9  & 4  & 5  & 44\% \\
\textsc{MathMutator}                & 4  & 0  & 4  & \textbf{0\%} \\
Alți mutatori                       & 74 & 44 & 33 & 59\% \\
\midrule
\textbf{Total}                      & \textbf{211} & \textbf{125} & \textbf{86} & \textbf{60\%} \\
\bottomrule
\end{tabularx}
\end{table}

\begin{figure}[!htb]
    \centering
    \includegraphics[width=\linewidth]{08-mutation-per-mutator.png}
    \caption{Distribuția mutanților pe tipuri (PITest). Verde = omorâți, roșu = supraviețuitori.}
    \label{fig:mutator-distribution}
\end{figure}

\subsection{Puncte forte ale suitei}

\begin{description}
    \item[\textsc{VoidMethodCall} - 100\% omorâți] Mutatorul șterge apelurile la metode \texttt{void}. Faptul că suita atinge 100\% înseamnă că toate apelurile la \texttt{repository.save(...)} sau \texttt{repository.delete(...)} sunt verificate de teste prin \texttt{verify(repo).save(...)} sau \texttt{verify(repo).delete(...)}. Dacă cineva ar șterge accidental aceste apeluri, am detecta imediat.

    \item[\textsc{NullReturnVals} - 94\% omorâți] Aproape toate metodele care returnează un obiect au testele care verifică non-null pe rezultat (cu \texttt{assertNotNull}). Doar 1 mutant a scăpat, ceea ce e excelent.
\end{description}

\subsection{Puncte slabe ale suitei}

\begin{description}
    \item[\textsc{MathMutator} - 0\% omorâți] Cei 4 mutanți pe operații aritmetice (\texttt{+/-/*//}) au scăpat toți. Asta indică o lacună în suită: nu am fost suficient de specifici cu valorile concrete în testele care verifică calcule. Vom adresa asta cu un test țintit (vezi secțiunea \ref{sec:killed-mutants}).

    \item[\textsc{ConditionalsBoundary} - 33\% omorâți] Doar 7 din 21 de mutanți pe granițe \texttt{>= / >} au fost omorâți. Asta indică faptul că nu testăm suficient de mult valori \textit{exact} pe granițe. Aici BVA-ul nostru nu e încă suficient de detaliat în toate metodele.
\end{description}

\section{Analiza unor mutanți supraviețuitori}

Pentru a înțelege concret unde sunt lacunele, am ales câțiva mutanți supraviețuitori reprezentativi din raport.

\subsection{Mutant supraviețuitor 1: \textsc{ConditionalsBoundary} pe linia 191}

Linia originală în \texttt{computeSummaryOwnedByUser}:

\begin{lstlisting}[language=Java]
if (percent >= 100) status = WalletSummary.Status.DEPASIT;
\end{lstlisting}

Mutantul a schimbat \texttt{>=} cu \texttt{>}. Pentru ca testul să detecteze diferența, ar fi nevoie ca:

\begin{itemize}
    \item un test cu \texttt{percent = 100} care \textbf{așteaptă} \texttt{DEPASIT};
    \item iar testul să verifice efectiv că \texttt{summary.getStatus() == DEPASIT}.
\end{itemize}

Avem un astfel de test în \texttt{...BoundaryTest.java} (\texttt{testSumar\_ExactPragDepasit\_100LaSuta}), iar în raportul PITest acest mutant apare \textit{killed}. Deci e omorât. \textbf{Vom vedea că un alt mutant similar, însă pe metoda \texttt{computeSmartBudgetAdvice}, e supraviețuitor} - ceea ce ne spune că trebuie să îl adresăm explicit.

\subsection{Mutant supraviețuitor 2: \textsc{MathMutator} pe linia 187}

Linia originală în \texttt{computeSummaryOwnedByUser}:

\begin{lstlisting}[language=Java]
int percent = totalSpent
    .multiply(BigDecimal.valueOf(100))
    .divide(budget, 0, RoundingMode.HALF_UP)
    .intValue();
\end{lstlisting}

Mutantul a schimbat \texttt{multiply} cu \texttt{divide} sau invers. Pentru ca testul să detecteze, trebuie să avem teste cu valori care produc rezultate \textit{exact} predictibile:

\begin{itemize}
    \item \texttt{totalSpent = 50}, \texttt{budget = 200} $\Rightarrow$ \texttt{percent = 25}.
    \item Dacă mutantul inversează operatorii, calculul ar produce o valoare cu totul diferită ($50/200 \neq 50 \cdot 100$).
\end{itemize}

Vom scrie un test țintit pentru asta în secțiunea următoare.

\section{Mutanții neechivalenți omorâți explicit}
\label{sec:killed-mutants}

Cerința T3 cere identificarea și omorârea a cel puțin doi mutanți \textit{neechivalenți}. Le-am ales pe cei doi cu cea mai mare valoare didactică.

\subsection{Mutantul 1: \textsc{MathMutator} pe \texttt{multiply} / \texttt{divide}}

\textbf{Linia originală:}

\begin{lstlisting}[language=Java]
int percent = totalSpent
    .multiply(BigDecimal.valueOf(100))
    .divide(budget, 0, RoundingMode.HALF_UP)
    .intValue();
\end{lstlisting}

\textbf{Mutația:} \texttt{multiply} schimbat cu \texttt{divide}, devenind:

\begin{lstlisting}[language=Java]
int percent = totalSpent
    .divide(BigDecimal.valueOf(100), 0, RoundingMode.HALF_UP)
    .multiply(budget)
    .intValue();
\end{lstlisting}

\textbf{De ce e neechivalent:} pentru \texttt{totalSpent = 50}, \texttt{budget = 200}, codul original returnează 25. Codul mutat (\texttt{50/100 * 200}) returnează 0 (după rotunjire), apoi \texttt{0 * 200 = 0}. Deci comportamentul e diferit observabil.

\textbf{Testul țintit:}

\begin{lstlisting}[language=Java, caption={Test țintit pentru mutantul \textsc{MathMutator}.}]
@Test
public void should_kill_percent_math_mutations() {
    wallet.setBudgetTotal(new BigDecimal("200"));
    when(txRepo.sumAmountByWalletId(any())).thenReturn(new BigDecimal("50"));

    WalletSummary s = service.computeSummaryOwnedByUser(tripId, login);

    // Daca multiply si divide sunt schimbate, calculul produce 0 in loc de 25
    assertEquals(25, s.getPercent());
}
\end{lstlisting}

\subsection{Mutantul 2: \textsc{ConditionalsBoundary} pe \texttt{percent >= 75}}

\textbf{Linia originală:} în \texttt{buildRiskUi}:

\begin{lstlisting}[language=Java]
if (percentUsed >= WARNING_THRESHOLD_PERCENT)
    return new RiskUi("Atenție", "bg-warning");
\end{lstlisting}

\textbf{Mutația:} \texttt{>=} schimbat cu \texttt{>}, devenind \texttt{percentUsed > 75}.

\textbf{De ce e neechivalent:} pentru \texttt{percent = 75}, codul original returnează \texttt{"Atenție"}. Codul mutat returnează \texttt{"OK"} (pentru că \texttt{75 > 75} e \texttt{false}).

\textbf{Testul țintit:}

\begin{lstlisting}[language=Java, caption={Test țintit pentru mutantul \textsc{ConditionalsBoundary}.}]
@Test
public void testSumar_ExactPragAtentie_75LaSuta() {
    wallet.setBudgetTotal(new BigDecimal("100"));
    when(txRepo.sumAmountByWalletId(any())).thenReturn(new BigDecimal("75"));

    WalletSummary s = service.computeSummaryOwnedByUser(tripId, login);

    assertEquals(75, s.getPercent());
    // Daca >= a fost mutat in >, statusul ar fi OK in loc de ATENTIE
    assertEquals(WalletSummary.Status.ATENTIE, s.getStatus());
}
\end{lstlisting}

\section{Cum identificăm mutanți echivalenți}

În raportul PITest, am întâlnit câțiva mutanți pe care i-am clasificat ca \textit{echivalenți} - adică alterații care nu schimbă comportamentul observabil al codului. Doi exemple:

\subsection{Exemplu 1: \texttt{++i} vs. \texttt{i++} într-un \texttt{for}}

În \texttt{computeInsightsOwnedByUser}, există un \texttt{for} clasic. PITest a generat un mutant care schimbă \texttt{i++} cu \texttt{++i}. În contextul unei bucle \texttt{for}, ambele variante au exact același efect (incrementarea se face după evaluarea condiției). Deci mutantul e echivalent prin definiție.

\subsection{Exemplu 2: schimbarea ordinii de evaluare la \texttt{a || b} fără side-effects}

Într-o condiție \texttt{a || b} unde nici \texttt{a} nici \texttt{b} nu au side-effects, schimbarea ordinii (\texttt{b || a}) nu afectează rezultatul. Dacă PITest generează un astfel de mutant pentru o condiție pură, e echivalent.

\section{Concluzii capitol}

Suita atinge un Mutation Score de \textbf{60\%}, peste pragul minim acceptabil (50\%) și aproape de țintele profesionale (75\%). Punctele forte ale suitei sunt verificarea side-effect-urilor pe repository (\textsc{VoidMethodCall} 100\%) și a returnurilor non-null (\textsc{NullReturnVals} 94\%).

Punctele slabe sunt pe operații aritmetice (\textsc{MathMutator}) și pe granițe (\textsc{ConditionalsBoundary}), pentru care am scris în mod explicit teste țintite. În total, am omorât 2 mutanți neechivalenți cu valoare didactică, conform cerinței T3.

În capitolul următor prezentăm experimentul cu un instrument AI (ChatGPT) și comparația suitei generate automat cu suita noastră.
